\subsection{Summary}
This work remarks how every innovation comes with perplexities, expecially when talking about security and privacy.
It was clear from the beginning that there is a lack of measures to ensure data privacy and control over useful and powerful tools as LLMs and agents in general.

The observation that started this thesis, that is having zero trust for both AI agent and its vendor, remain an open problem for enterpries an private customer alike.
Is important to state the vendors should provide more measures to handle agents and preferrably standardize those in a way more advanced works could be based on a more model-agnostic approach.

Ultimately, the hope of this work is to offer a solution to one important problem like the confidentiality over our own machine that now more than ever needs to be protected not only from what lies over the internet and other devices, but also from the potential threat that come from within ourself.
\subsection{Results}
\subsection{Future Development}\label{Future Development}
The implementation described in this work is a proof of concept and some of its choices do not reflect the deployment envisaged for a real organization.

Ideally, network traffic will travel inside a VPN and documents will be able to be decrypted not only from a client machine, but they may be choosen directly from a server, say, via an identifier.

As per today, the distance from the target deployment is summarized in the following table.
\begin{table}[htbp]
  \centering
  \begingroup\linespread{1.0}\small
  \begin{tabular}{@{}>{\raggedright\arraybackslash}m{0.17\textwidth}>{\raggedright\arraybackslash}m{0.36\textwidth}>{\raggedright\arraybackslash}m{0.39\textwidth}@{}}
    \toprule
    \textbf{Area} & \textbf{Prototype} & \textbf{Target deployment} \\
    \midrule
    Network & HTTP on a local network, ports of the services published to the host & every component inside a VPN, TLS between components, database not exposed \\
    Documents & read from the client machine and sent with each request; the server keeps no store & kept on the endpoint or retrieved from a central store by identifier \\
    Keycloak & development mode, data inside the container & production mode with a persistent, backed-up database \\
    Tokens & signature, issuer and expiry are verified & audience binding as well, as required by the MCP authorization specification \\
    Keys at rest & user-agent keys cached in PostgreSQL & keys protected at rest \\
    Client secrets & stored in a user directory of the client & stored in the secret store of the operating system \\
    Client enforcement & hook per user; managed deployment only for Claude Code & managed deployment for every harness that supports it \\
    Key cache & purged by hand if the authority is replaced & invalidated automatically when the authority changes \\
    Audit & JSON lines and a PostgreSQL mirror, local log not tamper-proof & append-only central storage \\
    \bottomrule
  \end{tabular}
  \endgroup
  \caption{From the prototype to the target deployment}\label{tab:target}
\end{table}

All the guarantees will hold and so will do the assumptions, all the future developments will increase the overall security and usability of this work.
